\documentclass[10pt,a4paper]{article}
\usepackage{webclases-lesson}
\lessonSetup{T0}{v0.4}{T0 - Herramientas para empezar Fisica}

\begin{document}
\shorthandoff{<>}
\color{Ink}

\coverHeader{T0 --- Herramientas para empezar Fisica}{Unidades, vectores, medida y graficas para resolver problemas con criterio.}

\begin{lessoncard}{Proposito}
Esta leccion es una caja de herramientas. No intenta convertir unidades, vectores, incertidumbres y calculo en cuatro temas completos: recoge lo minimo que conviene dominar antes de T1 y T2.
\begin{keybox}
\textbf{Clave de lectura.} Antes de calcular, escribe que magnitud buscas, en que unidades debe salir y que signo o direccion tendria sentido fisico.
\end{keybox}
\end{lessoncard}

\begin{lessoncard}{Indice por bloques}
\footnotesize
\begin{tabularx}{\linewidth}{@{}>{\raggedright\arraybackslash}p{1.6cm}Y>{\raggedleft\arraybackslash\bfseries}p{1.0cm}@{}}
\textbf{Bloque} & \textbf{Alcance} & \textbf{Pag.}\\
\midrule
\tocBlock{Bloque 1}{Magnitudes, unidades y comprobaciones dimensionales}
\tocLine{sec:b1}{1}{Medir, expresar unidades SI y detectar ecuaciones imposibles}
\tocBlock{Bloque 2}{Vectores, componentes y productos utiles}
\tocLine{sec:b2}{2}{Representar, proyectar, sumar y elegir el producto adecuado}
\tocBlock{Bloque 3}{Medida, incertidumbre y cifras significativas}
\tocLine{sec:b3}{3}{Distinguir resolucion, precision, error e incertidumbre}
\tocBlock{Bloque 4}{Graficas, derivadas e integrales en Fisica}
\tocLine{sec:b4}{4}{Leer pendientes, areas con signo y relaciones $x$-$v$-$a$}
\end{tabularx}
\end{lessoncard}

{\scriptsize\textcolor{Muted}{Fuentes: PDF original T0, background editorial T0 y contraste conceptual con OpenStax University Physics 1. Diagramas redibujados.}}

\clearpage

\lessonHeader{Bloque 1 \textperiodcentered\ Magnitudes, unidades y comprobaciones}{Apartados internos}{El lenguaje minimo del resultado}{Objetivo: que una ecuacion tenga sentido antes de sustituir numeros.}
\begin{lessoncard}{Magnitud, valor y unidad}
\subtopic{1}{Magnitudes, unidades y comprobaciones dimensionales}{sec:b1}
Una medida comunica tres cosas: la magnitud fisica, el valor numerico y la unidad. Escribir $12$ no basta; escribir $12\units{m}$ ya indica longitud. En problemas de Fisica, las unidades no decoran: controlan la coherencia del resultado.

\begin{center}
\begin{tabularx}{\linewidth}{@{}p{2.7cm}p{2.6cm}Y@{}}
\toprule
Tipo & Ejemplo & Uso rapido\\
\midrule
Fundamental SI & $m$, $kg$, $s$, $A$, $K$, $mol$, $cd$ & base de las demas\\
Derivada & $\N=kg\,m/s^2$, $\J=N\,m$ & comprueba formulas\\
Sin dimension & radian, coeficiente de rozamiento & puede aparecer en funciones o razones\\
\bottomrule
\end{tabularx}
\end{center}
\end{lessoncard}

\begin{examplebox}{analisis dimensional}
Se propone $v=v_0+2xt$ para una particula, donde $x$ es posicion y $t$ tiempo. Comprueba si puede ser una ecuacion de velocidad.
\solutionblock $[v]=L/T$ y $[v_0]=L/T$. En cambio $[2xt]=L\,T$. No se pueden sumar terminos con dimensiones distintas. La forma $v=v_0+at$ si es posible porque $[at]=(L/T^2)T=L/T$.
\variation{La comprobacion dimensional no demuestra que una formula sea verdadera; solo elimina expresiones imposibles.}
\end{examplebox}

\begin{warnbox}
Convertir unidades al final suele ocultar errores. Convierte al SI al principio salvo que el problema pida explicitamente otra unidad practica, como km/h.
\end{warnbox}

\clearpage

\lessonHeader{Bloque 1 \textperiodcentered\ Magnitudes, unidades y comprobaciones}{Aplicacion}{Convertir sin perder significado}{Objetivo: mantener el valor fisico aunque cambie la unidad.}
\begin{lessoncard}{Factores de conversion}
Un factor de conversion multiplica por $1$ escrito de otra forma. La magnitud no cambia; cambia la unidad usada para expresarla.
\[
72\units{km/h}\cdot\frac{1000\units{m}}{1\units{km}}\cdot\frac{1\units{h}}{3600\units{s}}=20,0\mps.
\]
Escribe cada factor para que la unidad que no quieres se cancele.
\end{lessoncard}

\begin{examplebox}{energia, trabajo y unidades}
Comprueba que $W=F\Delta x\cos\theta$ tiene unidades de energia.
\solutionblock
\[
[F\Delta x]=\N\,m=(kg\,m/s^2)m=kg\,m^2/s^2=\J.
\]
El coseno no tiene dimension. Si al sustituir obtienes $\N/s$ o $\kg/m$, el error esta antes del redondeo.
\end{examplebox}

\begin{lessoncard}{Mini-lista antes de calcular}
1. Convierte datos al SI. 2. Estima el orden de magnitud. 3. Escribe la unidad esperada. 4. Comprueba que todos los sumandos tienen las mismas dimensiones. 5. Al final, redondea segun la informacion disponible.
\end{lessoncard}

\clearpage

\lessonHeader{Bloque 2 \textperiodcentered\ Vectores, componentes y productos}{Apartados internos}{Direccion con algebra sencilla}{Objetivo: elegir ejes utiles y operar con signos visibles.}
\begin{lessoncard}{Vector y componentes}
\subtopic{2}{Vectores, componentes y productos utiles}{sec:b2}
Un vector queda definido por modulo, direccion y sentido. En ejes cartesianos:
\[
\vec A=A_x\ii+A_y\jj,\qquad |\vec A|=\sqrt{A_x^2+A_y^2},\qquad \tan\theta=\frac{A_y}{A_x}.
\]
\begin{center}
\begin{tikzpicture}[x=1cm,y=1cm]
  \draw[axis] (-.2,0) -- (5.2,0) node[right] {\scriptsize $x$};
  \draw[axis] (0,-.2) -- (0,3.2) node[above] {\scriptsize $y$};
  \draw[gridline] (3.8,0) -- (3.8,2.2) -- (0,2.2);
  \draw[vector] (0,0) -- (3.8,2.2) node[midway,above left] {\scriptsize $\vec A$};
  \draw[vectorgreen] (0,-.25) -- (3.8,-.25) node[midway,below] {\scriptsize $A_x=A\cos\theta$};
  \draw[vectoramber] (4.05,0) -- (4.05,2.2) node[midway,right] {\scriptsize $A_y=A\sin\theta$};
  \draw (0.75,0) arc(0:30:.75);
  \node at (.95,.25) {\scriptsize $\theta$};
\end{tikzpicture}
\end{center}
\end{lessoncard}

\begin{examplebox}{una fuerza inclinada}
$F=20,0\N$ forma $35^\circ$ sobre la horizontal. Halla componentes y decide que componente acelera un bloque sobre mesa lisa.
\solutionblock
\[
F_x=20,0\cos35^\circ=16,4\N,\qquad F_y=20,0\sin35^\circ=11,5\N.
\]
Si el bloque permanece en contacto con la mesa y no hay rozamiento, la aceleracion horizontal la produce $F_x$. La componente vertical modifica la normal: $N=mg-F_y$ si tira hacia arriba.
\end{examplebox}

\begin{lessoncard}{Productos que aparecen en Fisica}
\begin{tabularx}{\linewidth}{@{}p{3.0cm}Y Y@{}}
\toprule
Operacion & Resultado & Lectura fisica tipica\\
\midrule
$\vec A+\vec B$ & vector & resultante, desplazamiento total\\
$\lambda\vec A$ & vector & cambio de escala o sentido\\
$\vec A\cdot\vec B=AB\cos\theta$ & escalar & trabajo, proyeccion, alineacion\\
$|\vec A\times\vec B|=AB\sin\theta$ & vector perpendicular & momento, area orientada\\
\bottomrule
\end{tabularx}
\end{lessoncard}

\clearpage

\lessonHeader{Bloque 2 \textperiodcentered\ Vectores, componentes y productos}{Aplicacion}{Suma y proyeccion}{Objetivo: convertir dibujos en ecuaciones cortas.}
\begin{examplebox}{desplazamiento por tramos}
Un alumno camina $30\units{m}$ al este y despues $40\units{m}$ al norte.
\solutionblock
\[
\vec r=30\ii+40\jj,\qquad |\vec r|=50\units{m},\qquad \theta=\arctan(40/30)=53,1^\circ.
\]
La distancia recorrida es $70\units{m}$; el desplazamiento tiene modulo $50\units{m}$. No son la misma magnitud.
\end{examplebox}

\begin{center}
\begin{tikzpicture}[x=.9cm,y=.9cm]
  \draw[axis] (-.2,0)--(5,0) node[right] {\scriptsize este};
  \draw[axis] (0,-.2)--(0,4.5) node[above] {\scriptsize norte};
  \draw[vectorgreen] (0,0)--(3,0) node[midway,below] {\scriptsize $30$};
  \draw[vectorgreen] (3,0)--(3,4) node[midway,right] {\scriptsize $40$};
  \draw[vector] (0,0)--(3,4) node[midway,above left] {\scriptsize $\vec r$};
\end{tikzpicture}
\end{center}

\begin{lessoncard}{Como elegir ejes}
Elige ejes que hagan simples las fuerzas, velocidades o aceleraciones importantes. En un plano inclinado conviene un eje paralelo al plano; en proyectiles, horizontal y vertical; en circular, radial y tangencial.
\end{lessoncard}

\clearpage

\lessonHeader{Bloque 3 \textperiodcentered\ Medida e incertidumbre}{Apartados internos}{Cuanto confiar en un numero}{Objetivo: escribir resultados honestos y comparables.}
\begin{lessoncard}{Cuatro ideas que no son iguales}
\subtopic{3}{Medida, incertidumbre y cifras significativas}{sec:b3}
\begin{tabularx}{\linewidth}{@{}p{3.0cm}Y@{}}
\toprule
Idea & Significado practico\\
\midrule
Resolucion & cambio mas pequeno que distingue el instrumento\\
Precision & dispersion pequena al repetir medidas\\
Exactitud & cercania al valor de referencia\\
Error & diferencia respecto a un valor aceptado, si se conoce\\
Incertidumbre & intervalo razonable asociado al resultado comunicado\\
\bottomrule
\end{tabularx}
\end{lessoncard}

\begin{examplebox}{serie de medidas}
Se mide cinco veces una longitud en cm: $12,41$, $12,43$, $12,42$, $12,44$, $12,40$. El instrumento resuelve $0,01\units{cm}$.
\solutionblock Media:
\[
\bar L=\frac{12,41+12,43+12,42+12,44+12,40}{5}=12,42\units{cm}.
\]
La semidispersion es $(12,44-12,40)/2=0,02\units{cm}$. Un resultado razonable para este nivel es:
\[
L=(12,42\pm0,02)\units{cm}.
\]
No tiene sentido comunicar $12,420000\units{cm}$ si la incertidumbre esta en centesimas.
\end{examplebox}

\begin{exercisebox}{resultado con incertidumbre}
Un cronometro da $t=(2,36\pm0,04)\units{s}$ y una regla $d=(1,20\pm0,01)\units{m}$. Calcula la rapidez media y expresala con una incertidumbre razonable.
\solutionblock $v=d/t=0,508\mps$. Estimando incertidumbre relativa: $0,01/1,20+0,04/2,36\approx0,025$, luego $\Delta v\approx0,013\mps$. Resultado: $v=(0,51\pm0,01)\mps$.
\end{exercisebox}

\clearpage

\lessonHeader{Bloque 4 \textperiodcentered\ Graficas, derivadas e integrales}{Apartados internos}{Pendiente, area y movimiento}{Objetivo: leer informacion fisica en una grafica antes de usar formulas.}
\begin{lessoncard}{Pendiente como derivada}
\subtopic{4}{Graficas, derivadas e integrales en Fisica}{sec:b4}
Si $x(t)$ es posicion, la pendiente de la grafica posicion-tiempo es la velocidad. La pendiente media usa dos puntos; la instantanea usa la tangente.
\begin{center}
\begin{tikzpicture}[x=1.05cm,y=.75cm]
  \draw[axis] (0,0) -- (6.0,0) node[right] {\scriptsize $t$ (s)};
  \draw[axis] (0,0) -- (0,4.1) node[above] {\scriptsize $x$ (m)};
  \draw[BrandBlue,line width=1pt,domain=.2:5.4,samples=80] plot(\x,{0.25*(\x-1)^2+0.8});
  \draw[vectoramber] (2.0,1.05) -- (4.6,2.35) node[right] {\scriptsize pendiente media};
  \draw[vectorgreen] (3.2,1.95) -- (4.55,2.62) node[right] {\scriptsize tangente};
\end{tikzpicture}
\end{center}
\end{lessoncard}

\begin{examplebox}{de posicion a velocidad}
Para $x(t)=2+3t+t^2$ en SI, calcula $v(t)$ y la velocidad en $t=2,0\units{s}$.
\solutionblock
\[
v(t)=\frac{dx}{dt}=3+2t,\qquad v(2,0)=7,0\mps.
\]
La unidad confirma la lectura: derivar metros respecto a segundos da $\mathrm{m/s}$.
\end{examplebox}

\begin{lessoncard}{Area como integral}
El desplazamiento entre $t_1$ y $t_2$ se obtiene como area con signo bajo $v(t)$:
\[
\Delta x=\int_{t_1}^{t_2}v(t)\,dt.
\]
Area positiva aumenta $x$; area negativa disminuye $x$. La distancia recorrida suma areas en valor absoluto.
\begin{center}
\begin{tikzpicture}[x=1.0cm,y=.55cm]
  \draw[axis] (0,0) -- (6.1,0) node[right] {\scriptsize $t$ (s)};
  \draw[axis] (0,-2.0) -- (0,3.0) node[above] {\scriptsize $v$ (m/s)};
  \draw[fill=GreenLine!20,draw=none] (0,0) -- (2,2) -- (3,0) -- cycle;
  \draw[fill=AmberLine!20,draw=none] (3,0) -- (5,-1.4) -- (6,0) -- cycle;
  \draw[BrandBlue,line width=1pt] (0,0) -- (2,2) -- (3,0) -- (5,-1.4) -- (6,0);
  \node at (1.7,1.0) {\scriptsize area $+$};
  \node at (4.7,-.8) {\scriptsize area $-$};
\end{tikzpicture}
\end{center}
\end{lessoncard}

\begin{lessoncard}{Conexion final}
\[
a(t)=\frac{dv}{dt}=\frac{d^2x}{dt^2},\qquad
v(t)=v_0+\int_0^t a(\tau)\,d\tau,\qquad
x(t)=x_0+\int_0^t v(\tau)\,d\tau.
\]
En T1 usaremos estas ideas sin hacer calculo avanzado: pendiente de $x(t)$, pendiente de $v(t)$ y area bajo $v(t)$.
\end{lessoncard}

\end{document}
